\section{Markowitz utility of synthetic data}\label{sec:markowitz}
% =============================================================================================
% Owner: markowitz writer. Theory: docs/theory.md Section 9 (proofs in app:proofs-markowitz).
% Experiment: E4 (experiments/e4_markowitz.py), as run: 30 universes x 30 assets, B = 20
% replicas, 2 test seeds, 60 placebo replications per generator (results/e4/meta*.json).
% Main-text tables: e4_utility, e4_correlation (generated); manual_e4_tests_main,
% manual_e4_link_main (condensed, written by paper/scripts/e4_tests_main.py and
% e4_link_main.py from results/e4/*.csv). Full generated test and link tables: app:mk.
% Derived numbers not in a generated table (counts over universes, gross exposures, the sign
% count of the link tables, the 400 distinct tickers of results/e4/universes.csv, the reference
% start date 2020-07-24 = TRAIN day 643) were computed from results/e4/*.csv and the data
% loader; app:mk repeats them. Section 8.1 keeps the statements (lem:quad-loss ... ex:drift) that
% theory.tex (sec:theory-markowitz), appendix_proofs.tex and the introduction reference.
% Reliability table manual_e4_reliability (tab:manual-e4-reliability): seed-to-seed Spearman
% correlations of the four test quantities, written by paper/scripts/e4_reliability.py from
% results/e4/e4_tests.csv; its header comment records the critic's 60 to 63 training windows per
% sample (n_x - n_eval_x) and the 52097 trainable parameters of the E4 ConvCritic. The split of
% the 18/192 link count (rule: pooled view, universe-bootstrap interval excludes 0; within views,
% permutation p < 0.05; generator level, exact p < 0.05, never met; counting every printed interval
% instead would give 23) (4 positive: T vs pooled-moment GMV and GMV long-only, pooled and
% within-universe views; 14 negative: 8 z_hac, 6 wild-standardised MMD) was counted from
% results/e4/e4_link.csv. Zero risk-free rate: src/synthval/portfolio.py (rf = 0.0 default,
% not overridden in experiments/e4_markowitz.py).
% =============================================================================================

A two-sample test asks whether synthetic windows can be told apart from real ones; a user who
builds portfolios from synthetic data asks whether the portfolios perform well on real data. This
section asks whether the two answers agree. \Cref{sec:mk-theory} states what theory predicts,
\cref{sec:mk-design} describes the experiment E4, \cref{sec:mk-utility} reports the utility of
portfolios trained on synthetic data and \cref{sec:mk-link} relates it to multivariate \WC{} and
MMD tests. The answer is negative: the generator that the tests reject most often loses almost
nothing for the minimum-variance portfolio, the generator with the largest losses is rejected,
after size adjustment, less often than the two \iid{} generators, and the one quantity that
tracks the minimum-variance loss is the distance between the synthetic and the past average
correlation, with a negative sign, because the real correlation fell between \TRAIN{} and
\TEST{}. Three qualifications apply. The evidence comes from one evaluation period, in which
correlations fell, and carries no inference over dates; the portfolios use each generator only
through the unconditional moments of long simulated samples; and within a generator the
multivariate critic statistic is too noisy across synthetic seeds to show clearly whether it
tracks the loss.

% ---------------------------------------------------------------------------------------------
\subsection{What theory predicts}\label{sec:mk-theory}

\Cref{sec:theory-markowitz} summarised what the theory implies for E4; this subsection states the
results. Let $r_t\in\R^N$ be the daily returns of $N$ assets, with real \TRAIN{} sample $R_{\mathrm{tr}}$
($\Ttr\times N$) and real \TEST{} sample $R_{\mathrm{te}}$ ($\Tte\times N$), and let $(\mu,\Sigma)$
be the \TEST{} moments: population moments, or the realised moments
$(\hat\mu_{\mathrm{te}},S_{\mathrm{te}})$, for which every statement below is an identity for the
measured quantities. Portfolios lie in $\mathcal W=\{w\in\R^N:\one'w=1\}$, intersected with
$\{w\ge0\}$ for long-only rules. Given estimates $(\hat\mu,\hat\Sigma)$, the global
minimum-variance (GMV) portfolio is $\wgmv(\hat\Sigma)=\argmin_{w\in\mathcal W}w'\hat\Sigma w$
($=\hat\Sigma^{-1}\one/(\one'\hat\Sigma^{-1}\one)$ without the sign constraint) and the mean-variance
portfolio with risk aversion $\gamma$ is
$\wmv(\hat\mu,\hat\Sigma)=\argmax_{w\in\mathcal W}\{w'\hat\mu-\tfrac\gamma2w'\hat\Sigma w\}$
\citep{markowitz1952portfolio}. Trained on real data (TRTR), these maps are applied to the moments
of $R_{\mathrm{tr}}$ and give $\hat w_{\mathrm{real}}$; trained on synthetic data (TSTR), to the
moments of a synthetic sample $Y^{(b)}$ of length $\Ttr$ from a generator $G$ calibrated on
$R_{\mathrm{tr}}$, and give $\hat w^{(b)}_{\mathrm{syn}}$, $b=1,\dots,B$. With $V(w)=w'\Sigma w$,
$w^*=\wgmv(\Sigma)$ and $V^*=V(w^*)$, the losses are
\begin{gather}
  \mathcal L_{\mathrm{gmv}}(w)=V(w)-V^*,\label{eq:gmv-loss}\\
  \mathcal L_{\mathrm{mv}}(w)=\CE(w^*_{\mathrm{mv}})-\CE(w),\qquad
  \CE(w)=w'\mu-\tfrac\gamma2\,w'\Sigma w,\qquad w^*_{\mathrm{mv}}=\wmv(\mu,\Sigma).\label{eq:ce}
\end{gather}
E4 reports the variance ratio $V(\hat w_{\mathrm{syn}})/V(\hat w_{\mathrm{real}})$, which exceeds
one exactly when $\mathcal L_{\mathrm{gmv}}(\hat w_{\mathrm{syn}})>\mathcal L_{\mathrm{gmv}}(\hat w_{\mathrm{real}})$,
and the certainty-equivalent (CE) loss
$\CE(\hat w_{\mathrm{real}})-\CE(\hat w_{\mathrm{syn}})=\mathcal L_{\mathrm{mv}}(\hat w_{\mathrm{syn}})-\mathcal L_{\mathrm{mv}}(\hat w_{\mathrm{real}})$.
Proofs are in \cref{app:proofs-markowitz}.

\begin{lemma}[Quadratic loss identities]\label{lem:quad-loss}
(a) If $\mathcal W=\{w:\one'w=1\}$, then $V(w)-V^*=(w-w^*)'\Sigma(w-w^*)$ for every $w\in\mathcal W$.
(b) If $\mathcal W$ is closed and convex and $w^*$ minimises $V$ over $\mathcal W$, then
$V(w)-V^*\ge(w-w^*)'\Sigma(w-w^*)$ for every $w\in\mathcal W$.
(c) If $\mathcal W=\{w:\one'w=1\}$, then
$\mathcal L_{\mathrm{mv}}(w)=\tfrac\gamma2(w-w^*_{\mathrm{mv}})'\Sigma(w-w^*_{\mathrm{mv}})$; under
further convex constraints the equality becomes ``$\ge$''.
\end{lemma}

With $\Sigma=S_{\mathrm{te}}$, part (a) decomposes the measured GMV gap exactly into a difference
of squared Mahalanobis distances to the ex-post GMV portfolio,
$\norm{\hat w_{\mathrm{syn}}-w^*}_{S_{\mathrm{te}}}^2-\norm{\hat w_{\mathrm{real}}-w^*}_{S_{\mathrm{te}}}^2$,
where $\norm{x}_S^2=x'Sx$.

\begin{proposition}[Sensitivity of the GMV portfolio to the covariance]\label{prop:gmv-sensitivity}
Let $\mathcal W=\{w:\one'w=1\}$, let $\Sigma$ and $\hat\Sigma$ be symmetric positive definite,
$A=\Sigma^{-1/2}\hat\Sigma\Sigma^{-1/2}$, $\kappa=\lambda_{\max}(A)/\lambda_{\min}(A)=\cond(\Sigma^{-1}\hat\Sigma)$
and $M=\Sigma^{-1}-\Sigma^{-1}\one\one'\Sigma^{-1}/(\one'\Sigma^{-1}\one)$, and let
$\norm{\cdot}_{\mathrm{op}}$ be the spectral norm.
\begin{enumerate}
  \item[(a)] (Scale invariance.) $\wgmv(c\hat\Sigma)=\wgmv(\hat\Sigma)$ for every $c>0$; in
  particular $V(\wgmv(c\Sigma))=V^*$.
  \item[(b)] (Kantorovich bound.) $1\le V(\wgmv(\hat\Sigma))/V^*\le(1+\kappa)^2/(4\kappa)$, and the
  upper bound is sharp.
  \item[(c)] (Quadratic in the relative error.) If $\varepsilon_0=\norm{\Sigma^{-1/2}(\hat\Sigma-\Sigma)\Sigma^{-1/2}}_{\mathrm{op}}$
  is below one, then $V(\wgmv(\hat\Sigma))/V^*-1\le\varepsilon_0^2/(1-\varepsilon_0^2)$.
  \item[(d)] (First order.) As $E\to0$, $\wgmv(\Sigma+E)-w^*=-MEw^*+O(\norm{E}^2)$ and
  $V(\wgmv(\Sigma+E))-V^*=(Ew^*)'M(Ew^*)+O(\norm{E}^3)$, where $M\succeq0$ has kernel spanned by
  $\one$: the quadratic term vanishes if and only if $w^*$ is also the GMV portfolio of $\Sigma+E$.
\end{enumerate}
\end{proposition}

\begin{corollary}[What drives the GMV utility of a generator]\label{cor:gmv-drivers}
Let $\hat\Sigma_{\mathrm{syn}}$ be the sample covariance of a synthetic sample of length $T_s$ from
a stationary ergodic generator with covariance $\Sigma_G$, and $\kappa_G=\cond(\Sigma^{-1}\Sigma_G)$.
\begin{enumerate}
  \item[(i)] As $T_s\to\infty$,
  $V(\wgmv(\hat\Sigma_{\mathrm{syn}}))/V^*\to V(\wgmv(\Sigma_G))/V^*\le(1+\kappa_G)^2/(4\kappa_G)$
  almost surely: in the large-sample limit the GMV utility of a generator depends on it only
  through $\Sigma_G$, and only up to scale.
  \item[(ii)] (Heuristic.) At finite $T_s$,
  $\E\,\mathcal L_{\mathrm{gmv}}\approx\{(\Sigma_G-\Sigma)w^*\}'M\{(\Sigma_G-\Sigma)w^*\}+T_s^{-1}\Tr(M\LRV_G)$,
  with $\LRV_G$ the long-run covariance of $r_t(r_t'w^*)$ under $G$: tails and volatility
  clustering enter only through the second, estimation-noise term.
  \item[(iii)] (Heuristic.) If $\Sigma_G\approx\hat\Sigma_{\mathrm{tr}}$, the \TRAIN{} sample
  covariance, and $T_s=\Ttr$, then
  $\hat\Sigma_{\mathrm{syn}}-\Sigma=(\hat\Sigma_{\mathrm{tr}}-\Sigma)+(\hat\Sigma_{\mathrm{syn}}-\Sigma_G)$
  with two conditionally uncorrelated terms of comparable size: one synthetic sample roughly
  doubles the estimation-noise part of the TRTR loss, and pooling $B$ samples divides the added
  part by $B$. The same holds, with more force, for the mean in mean-variance portfolios.
\end{enumerate}
\end{corollary}

\begin{proposition}[No free lunch from synthetic data]\label{prop:no-free-lunch}
Let $Y=g(R_{\mathrm{tr}},U)$ with $U$ independent of $(R_{\mathrm{tr}},\Theta)$, where $\Theta$ is
any \TEST{} quantity (such as $\Sigma$, $(\mu,\Sigma)$ or $R_{\mathrm{te}}$). Let $\pi$ be a
portfolio rule with values in a convex set $\mathcal W$ and $\ell(w;\Theta)$ be convex in $w$.
Then $\bar w=\E[\pi(Y)\mid R_{\mathrm{tr}}]$ lies in $\mathcal W$, uses the data only through
$R_{\mathrm{tr}}$, and $\ell(\bar w;\Theta)\le\E[\ell(\pi(Y);\Theta)\mid R_{\mathrm{tr}},\Theta]$.
\end{proposition}

Since $V$ and $-\CE$ are convex, no TSTR rule beats, on average over the synthetic randomness,
the best rule based on the real past alone, and averaging the weights over replicas does at least
as well as averaging the losses of replica-wise portfolios. Synthetic data can beat a
\emph{particular} TRTR rule, such as the sample-covariance GMV portfolio, only by acting as a
regulariser, as covariance shrinkage does \citep{ledoit2004well,jagannathan2003risk}. Test
distance and utility can disagree in both directions:

\begin{example}[Rejected by the critic, perfect for GMV]\label{ex:scale-blind}
Let returns be \iid{} $N(0,\Sigma)$ and $G$ \iid{} $N(0,c\Sigma)$ with $c\neq1$. Windows satisfy
$Y\eqd\sqrt c\,X$, so $\W(P,Q)=\abs{\sqrt c-1}\,\E\norm{X}>0$ and a consistent critic test rejects
with probability tending to one (\cref{prop:consistency}); by \cref{prop:gmv-sensitivity}(a) the
GMV loss of $\Sigma_G=c\Sigma$ is zero.
\end{example}

\begin{example}[Nearly invisible in $\W$, large GMV loss]\label{ex:spread}
Let $N=2$, $v_1=(\cos\phi,\sin\phi)'$, $v_2=(-\sin\phi,\cos\phi)'$, $\Sigma=v_1v_1'+\eta^2v_2v_2'$
and $\Sigma_G=v_1v_1'+s\eta^2v_2v_2'$ with $0<s,\eta<1$ and
$\tan\phi=(1-\eta\sqrt s)/(1+\eta\sqrt s)$. Then $\kappa_G=1/s$ and
$V(\wgmv(\Sigma_G))/V^*=(1+s)^2/(4s)$, the Kantorovich bound, which equals $1.5625$ for $s=1/4$
whatever $\eta$. For \iid{} Gaussian windows of length $L$, the quadratic Wasserstein distance
gives $\W\le W_2=\sqrt L\,\eta(1-\sqrt s)\to0$ as $\eta\to0$. The error lies in the low-variance
spread direction $v_2$, which GMV weights heavily (for $\eta=0.1$, $w^*=(-1.48,2.48)'$ and
$\wgmv(\Sigma_G)=(-4.49,5.49)'$) and which $\W$ discounts by $\sqrt{\lambda_2}=\eta$.
\end{example}

\begin{example}[Univariate tests cannot see correlation]\label{ex:corr-blind}
A generator with exactly the univariate window laws of the data and independent assets satisfies
every per-ticker null of E3, so each univariate test rejects it with probability at most $\alpha$;
its covariance is $\diag(\Sigma)$, and its GMV portfolio can be arbitrarily bad.
\end{example}

\begin{example}[Drift errors are undetectable yet decisive]\label{ex:drift}
Let returns be \iid{} Gaussian with known standard deviation $\sigma$, and let the synthetic drift
of one asset be off by $\delta_\mu$. By \cref{prop:oracle-bound}(b), no level-$\alpha$ two-sample
test that uses $\Tte$ real days has power above
$1-\Phi(z_{1-\alpha}-\abs{\delta_\mu}\sqrt{\Tte}/\sigma)$, however many synthetic draws it uses.
For $\abs{\delta_\mu}=\sigma/\sqrt{\Ttr}$, one standard error of the
\TRAIN{} mean, and $\Tte/\Ttr=615/1258$, the bound is $1-\Phi(1.645-0.70)\approx0.17$ at
$\alpha=0.05$. Errors of this size dominate mean-variance utility
\citep{michaud1989markowitz,chopra1993effect}: with $\hat\Sigma=\Sigma$ and the budget constraint
only, \cref{lem:quad-loss}(c) gives the loss $(2\gamma)^{-1}(\hat\mu-\mu)'M(\hat\mu-\mu)$.
\end{example}

For E4 this predicts that GMV losses track the accuracy of the covariance up to scale, not the
rejection rate; that mean-variance losses are dominated by mean estimation noise, which no test
detects and pooling replicas reduces; and that the covariance that matters is that of \TEST{}, so
that fidelity to the past and usefulness for the future differ when the covariance changes.

% ---------------------------------------------------------------------------------------------
\subsection{Design of E4}\label{sec:mk-design}

\paragraph{Universes and generators.}
E4 draws 30 universes of $N=30$ assets, each without replacement and with its own seed, from the
479 names of the complete universe (\cref{sec:data-panel}); universes may share tickers, and
together they contain 400 distinct names. Five
multivariate generators are fitted to the real \TRAIN{} returns of each universe: MV Gaussian
(\iid{}, sample mean and covariance); MV Student-$t$ (\iid{}, location, scatter and degrees of
freedom by ECME; \citealp{liu1995ml}); CCC-GARCH-$t$ (\GARCHt{} margins with variance targeting
and persistence cap 0.995, Gaussian copula with constant correlation estimated from the normal
scores of the standardised residuals; \citealp{bollerslev1990modelling,joe2005asymptotic}); MV
bootstrap (stationary bootstrap of whole \TRAIN{} days with blocks shared by all assets;
\citealp{politis1994stationary}); and factor GARCH-$t$ (single index equal to the equal-weight
return of the same 30 assets, OLS loadings, \GARCHt{} factor and independent \GARCHt{} residuals).
Fitting details are in \cref{app:impl} and diagnostics in \cref{app:mk-generators}.

\paragraph{Utility.}
For each universe and generator, $B=20$ synthetic samples of $\Ttr=1258$ days are simulated from a
stationary start (burn-in 250 days). The portfolios therefore use a generator only through the
moments of long simulated samples, essentially its unconditional mean and covariance: the
conditional covariance forecasts that CCC-GARCH-$t$ and factor GARCH-$t$ could deliver from their
end-of-\TRAIN{} state are deliberately not used, and a variant that would use them (for instance
\TEST{}-length samples started from that state, as in the tests below) was not run. Comparisons
between dynamic and \iid{} generators below concern their unconditional covariances, not their
forecasts. On each sample, and on $R_{\mathrm{tr}}$, the mean is the sample mean and the
covariance the sample covariance or the Ledoit--Wolf estimator with identity target
\citep{ledoit2004well}. Six rules are built: GMV, GMV long-only, long-only mean-variance with
$\gamma\in\{1,5,20\}$ and $1/N$ \citep{demiguel2009optimal}. The weights are held fixed with daily
rebalancing and evaluated on the 615 \TEST{} days through daily simple returns, against the same
rule trained on $R_{\mathrm{tr}}$. The losses are the variance ratio, the Sharpe loss
$\SR(\hat w_{\mathrm{real}})-\SR(\hat w_{\mathrm{syn}})$ (annualised, with a zero risk-free
rate) and the CE loss with the $\gamma$ of the rule ($\gamma=5$ for GMV), in percent per
year ($252\times100$ times the daily value); ratios above one and positive losses mean that TSTR
is worse. Every portfolio is fully invested, so a positive risk-free rate would leave the weights,
the variance ratios and the CE losses unchanged; it would lower the Sharpe ratios and change the
Sharpe losses between portfolios of different volatility. The returns are price returns
(\cref{sec:data-panel}), so the \TRAIN{} means and the levels of the certainty equivalents and
Sharpe ratios omit the dividend yield, and the selection of the universe on membership at the end
of \TEST{} biases the \TRAIN{} means upward. The losses are aggregated over
replicas in three ways: \emph{replica-wise}, the mean loss of the $B$ single-replica portfolios
(the DESIGN measure, which contains the added noise of \cref{cor:gmv-drivers}(iii));
\emph{pooled}, the loss of one portfolio built from the moments of the $B$ stacked replicas, which
isolates misspecification; and \emph{averaged weights} (\cref{prop:no-free-lunch}). The $1/N$
gaps are zero by construction. No inference over dates is provided. All losses are point
estimates on the same 615 \TEST{} days, which every universe shares, so a statement that holds in
all 30 universes describes one realisation seen through 30 overlapping portfolios, not 30
independent confirmations; a test over \TEST{} days, such as the HAC test of Sharpe-ratio
differences of \citet{ledoit2008robust} or its analogue for variances, was not run.

\paragraph{Tests.}
For each universe and generator, one synthetic sample of \TEST{} length is simulated from the end
of \TRAIN{} (GARCH filters continue from their \TRAIN{} state). It and the real \TEST{} sample are
divided by the per-ticker \TRAIN{} standard deviation and cut into non-overlapping windows of
shape $(N,L)=(30,5)$, 123 per sample. The tested object thus differs from the one used for
utility: one \TEST{}-length sample started from the end-of-\TRAIN{} state, against $B$
\TRAIN{}-length samples from a stationary start; for the \iid{} generators the two differ only in
length, for the others also in the starting state. The \WC{} test uses the convolutional critic of
\cref{sec:method-critic}, the time-block split and $\phac$. Each critic is trained on the 60 to 63
windows per sample of its training half, each of dimension $30\times5=150$, far fewer than the
52\,097 trainable parameters of the network. No architecture or training ablation was run for
this multivariate critic (the ablation of \cref{sec:robust} concerns the univariate critic), and
\cref{sec:mk-link} shows that its statistic is noisy across synthetic seeds. MMD uses the wild
bootstrap with a bandwidth fixed per universe (median heuristic of the real \TEST{} windows), so
that all generators of a universe face the same kernel, and they also share the split, critic
initialisation and mini-batches (the critic's input standardisation still uses the pooled
training half, which depends on the generator). Two synthetic seeds were run (DESIGN: ten;
compute budget); rejection shares across universes are reported for the first seed and for $\pmed$
\eqref{eq:pmed}. A reference row uses the last 615 days of real \TRAIN{} (2020-07-24 to
2022-12-30) in place of a synthetic sample. Because $\phac$ and $\pwild$ need not be calibrated
with 123 windows and persistent volatility, 60 placebo replications per generator (DESIGN: 200)
apply the same call to two independent draws of the fitted generator; they give the size and
size-adjusted Monte Carlo p-values (\cref{app:mk-placebo}). E4 uses the absolute tests only: no
anchored relative test, power check or verdict is computed for the multivariate generators.

\paragraph{Link.}
Each (universe, generator) pair is one point, with the seed means of $T$, $\zhac$, $\TMMD$ and the
wild-standardised MMD statistic, and a benchmark that is not a test: the absolute difference
$\abs{\Delta\,\text{avg.\ corr.}}$ between the mean pairwise correlation of the synthetic replicas
and that of real \TRAIN{}. Spearman correlations with the losses are computed pooled over the 150
points, within universes (across the five generators), within generators (across the 30
universes, which removes generator-level offsets) and between the five generator means
(\cref{app:mk-link}). Universes share every \TEST{} date and overlap in tickers, so intervals that
treat them as independent are optimistic; conclusions are conditional on the five generators.
The reproducibility of each test quantity is measured by the Spearman correlation between its
values for the two synthetic seeds (\cref{tab:manual-e4-reliability}).

% ---------------------------------------------------------------------------------------------
\subsection{Utility of portfolios trained on synthetic data}\label{sec:mk-utility}

\begin{table}[tbp]
  \centering
  \caption{Utility losses of portfolios trained on synthetic data relative to the same rule
  trained on real \TRAIN{} (E4). Means over 30 universes of 30 assets of the replica-wise losses
  ($B=20$ synthetic \TRAIN{}-length samples per universe and generator), evaluated on the 615
  \TEST{} days. Var.\ ratio: \TEST{} variance of the TSTR portfolio over that of the TRTR
  portfolio; Sharpe loss: annualised Sharpe ratio of TRTR minus TSTR; CE loss: certainty
  equivalent of TRTR minus TSTR, with the $\gamma$ of the rule ($\gamma=5$ for GMV), percent per
  year. Ratios above one and positive losses: TSTR worse. Sample: sample covariance; LW:
  Ledoit--Wolf (identity target), used for both TSTR and TRTR. Point estimates; no intervals.
  Other aggregations: \cref{tab:e4-utility-agg}.}
  \label{tab:e4-utility}
  \small
  \fittable{e4_utility}
\end{table}

\begin{table}[tbp]
  \centering
  \caption{Correlation structure of real and synthetic samples (E4). Avg.\ corr.: mean pairwise
  return correlation, mean over the 30 universes (sd across universes); synthetic values are
  means over the 20 \TRAIN{}-length replicas. Log-vol corr.: mean pairwise correlation of
  rolling 20-day log-volatilities. $\kappa$: median over universes of
  $\cond(S_{\mathrm{te}}^{-1}\hat\Sigma)$, with $\hat\Sigma$ the real \TRAIN{} covariance or that
  of a synthetic replica (median over replicas first); Kant.\ bound: the corresponding median of
  $(1+\kappa)^2/(4\kappa)$, the bound of \cref{prop:gmv-sensitivity}(b). At cap: GARCH components
  whose persistence reached the cap 0.995, out of all components (CCC-GARCH-$t$: 30 margins per
  universe; factor GARCH-$t$: factor and 30 residuals).}
  \label{tab:e4-correlation}
  \small
  \fittable{e4_correlation}
\end{table}

\Cref{tab:e4-utility} reports the replica-wise losses and \cref{tab:e4-correlation} the
correlation structure; the TRTR benchmark and the other aggregations are in
\cref{app:mk-aggregation}. Ledoit--Wolf moments change the losses little; we discuss sample
moments.

\paragraph{Minimum-variance portfolios.}
Four generators give GMV portfolios whose \TEST{} variance is within 6\% of the TRTR variance
(replica-wise ratios 0.954 for MV Student-$t$, 1.021 for CCC-GARCH-$t$, 1.038 for MV Gaussian,
1.054 for the MV bootstrap); factor GARCH-$t$ has 1.436. With pooled moments the ratios are 0.896,
0.869, 1.004, 0.993 and 1.466 in the same order (\cref{tab:e4-utility-agg}): pooling reduces the
noise but does not help factor GARCH-$t$, whose loss is misspecification. With pooled moments,
CCC-GARCH-$t$ and MV Student-$t$ build better GMV portfolios than the real data they were fitted to
in all 30 universes, and factor GARCH-$t$ has the largest replica-wise ratio of the five generators
in all 30 (counts from the E4 output files; as noted in \cref{sec:mk-design}, these are 30 views
of one \TEST{} period).

\paragraph{The correlation regime shift.}
The mean pairwise correlation fell from 0.419 in \TRAIN{} to 0.268 in \TEST{}, in each of the 30
universes, and the correlation of log-volatilities from 0.603 to 0.302 (\cref{tab:e4-correlation});
\TRAIN{} contains the market-wide crash of 2020. Generators calibrated on \TRAIN{} inherit
the past correlation to different degrees: MV Gaussian exactly (0.419), the MV bootstrap and factor
GARCH-$t$ nearly (0.406, 0.398), MV Student-$t$ less (0.359; plausibly because the $t$ scatter
estimate downweights days with a large Mahalanobis distance, among them the market-wide days of
2020) and CCC-GARCH-$t$ least (0.310), a documented attenuation of the model, whose volatilities
co-move only through the correlation of the shocks (log-volatility correlation 0.127). By
fidelity to the past CCC-GARCH-$t$ is the worst generator; by usefulness for the future, measured
by the pooled-moment GMV variance, it is the best. Lower correlations act like shrinkage of the
covariance towards its diagonal, which is the direction in which the real covariance moved: its
GMV weights are less leveraged (mean gross exposure $\sum_i\abs{w_i}$ of the replica-wise
portfolios 1.706 against 1.986 for TRTR), and its pooled-moment variance ratio is the lowest in 19
of the 30 universes (MV Student-$t$ in 10). This does not contradict \cref{prop:no-free-lunch}:
the attenuated covariance is a regulariser that happened to point in the direction of the regime
change, an advantage that a different future could reverse. Whether the real past alone, treated
in the same spirit, achieves the same gain was not examined: shrinkage towards a
constant-correlation target \citep{ledoit_2004_honey}, a covariance estimated on the most recent
part of \TRAIN{} (such as its last 615 days, the reference row of \cref{tab:e4-tests}) and EWMA or
DCC forecasts \citep{engle2002dynamic} were not run. They would be needed to show that synthetic
data add anything beyond a shrinkage of the correlations.

Factor GARCH-$t$ illustrates the mechanism of \cref{ex:spread}. It matches the average correlation
and the co-movement of volatilities (0.374) better than CCC-GARCH-$t$, yet it has the largest GMV
loss and the largest condition number (31.16, Kantorovich bound 8.299). Two features of its
construction plausibly contribute, and E4 does not separate them (\cref{app:mk-generators}).
First, its covariance is one factor plus a diagonal, so the residual covariances, which carry the
industry structure, are zero. Second, the factor is the equal-weight return of the same 30
assets, so in the data the OLS residuals add up to zero every day and have negative covariances
on average; the simulated residuals are independent, violate this identity and add idiosyncratic
variance to the equal-weight direction. Both errors change the low-variance spread directions on
which GMV portfolios load, while the average correlation is right. An external factor, such as
the equal-weight return of all 479 names or an index, or factors from principal components would
remove the adding-up artefact and leave only the omitted industry structure; this variant was not
run. The Kantorovich bound held in every replica but does not rank the generators: CCC-GARCH-$t$
has a larger $\kappa$ (24.38) than the real \TRAIN{} covariance (18.65) and the lowest pooled
variance ratio.

\paragraph{Mean-variance portfolios.}
For every generator and mean-variance rule the replica-wise TSTR portfolio has a larger \TEST{}
variance and a lower CE than the TRTR portfolio: variance ratios from 1.117 to 1.597, CE losses
from 0.32 to 5.03 percent per year. These losses are mostly noise in the synthetic means: pooling
the replicas lowers the CE loss in 14 of the 15 pairs of generator and rule (for MV Gaussian with
$\gamma=1$ from 3.57 to $-1.36$), as \cref{cor:gmv-drivers}(iii) predicts, and averaged weights do
at least as well as the replica-wise portfolios in every universe, generator and rule, as
\cref{prop:no-free-lunch} requires. Even trained on real data, no mean-variance rule beats $1/N$
in Sharpe ratio (0.72 for $\gamma=5$ against 0.83, with a zero risk-free rate;
\cref{tab:e4-trtr}), as in \citet{demiguel2009optimal}.

% ---------------------------------------------------------------------------------------------
\subsection{Do test results predict utility?}\label{sec:mk-link}

\begin{table}[tbp]
  \centering
  \caption{Rejection shares of the multivariate tests across the 30 universes at $\alpha=0.05$
  (E4): real \TEST{} against one synthetic \TEST{}-length sample per seed, windows
  $(N,L)=(30,5)$, 123 per sample. \WC{}: time-block split, Driscoll--Kraay p-value; MMD: wild
  bootstrap, bandwidth fixed per universe. First: first synthetic seed; $\pmed$:
  $\min\{1,2\med_sp_s\}$ over two seeds. Placebo size: rejection rate when both samples are
  independent draws from the fitted generator (60 replications per generator,
  \cref{tab:e4-placebo}). Size-adjusted: p-values calibrated against the placebo scores of the
  same generator and test. Reference row: the last 615 days of real \TRAIN{} instead of a
  synthetic sample; it has no placebo. Rates only; Wilson 95\% intervals, which treat universes as
  independent, are in \cref{tab:e4-tests-full}.}
  \label{tab:e4-tests}
  \small
  \fittable{manual_e4_tests_main}
\end{table}

\paragraph{What the tests say.}
With nominal p-values and the first seed, the \WC{} test rejects each generator in 0.63 to 1.00
of the universes and MMD in 0.33 to 0.97, and the \WC{} test rejects the recent real past in 0.67
of them (\cref{tab:e4-tests}). The placebo shows that the
nominal $\phac$ is oversized for generators with persistent volatility: 0.27 at 5\% for
CCC-GARCH-$t$ and factor GARCH-$t$ and 0.13 for the MV bootstrap, against 0.05 and 0.07 for the
\iid{} generators (\cref{fig:e4-placebo-size}; the cause is discussed in \cref{app:mk-placebo}).
The reference row cannot be size-adjusted. After size adjustment the \WC{} test still rejects MV
Gaussian in all universes and MV Student-$t$ in 0.87 of them (first seed), but CCC-GARCH-$t$, the
MV bootstrap and factor GARCH-$t$ only in 0.33, 0.33 and 0.40 (0.07, 0.10 and 0.27 with $\pmed$).
MMD also separates the \iid{} generators from the dynamic ones and rejects the latter less often.
Both tests find the generators without volatility clustering easiest to distinguish from real
\TEST{} windows.

\paragraph{What the tests do not say.}
This ordering is unrelated to utility. MV Gaussian is rejected in every universe after size
adjustment (first seed) and has the largest critic statistic of the five generators in 17 of the
30 universes, yet its replica-wise GMV variance ratio lies between 1.007 and 1.060 in every
universe and its pooled-moment ratio is 1.004 on average. No witness analysis was run for the
multivariate critic, so what it detects is not known. If it behaves like the univariate critic of
\cref{sec:emp-witness}, which is close to a realised-volatility statistic, it responds mainly to
the level of window volatility and, for a generator without clustering, to the missing dispersion
of that level across windows. A common volatility level acts as a scale, to which the GMV
portfolio is blind (\cref{ex:scale-blind}), and the dispersion enters the GMV loss only through
estimation noise (\cref{cor:gmv-drivers}(ii)). Factor GARCH-$t$ is rejected less often than
either \iid{} generator after size adjustment, and has the largest GMV loss in every universe.
Mean-variance losses come from mean errors that no test can see (\cref{ex:drift}).

\begin{table}[tbp]
  \centering
  \caption{Reproducibility of the E4 test quantities across synthetic seeds: Spearman
  correlation between the values of the first and the second synthetic seed. Within generator:
  across the 30 universes of one generator (the variation used by the within-generator view of
  \cref{tab:e4-link}), and the mean of the five values. Pooled: over the 150 (universe,
  generator) points. Reference: last 615 days of real \TRAIN{}, the same data in both seeds, so
  that only the randomness of the test differs (split, initialisation and mini-batches of the
  critic; wild-bootstrap multipliers of MMD). \WC{}: time-block split, critic trained on 60 to 63
  windows of dimension 150 per sample; MMD: bandwidth fixed per universe.}
  \label{tab:manual-e4-reliability}
  \small
  \fittable{manual_e4_reliability}
\end{table}

\paragraph{How reproducible the test quantities are.}
A correlation between a test quantity and a loss across universes is informative only if the
quantity itself is reproducible across universes. \Cref{tab:manual-e4-reliability} shows that,
within a generator, the critic statistic is not. Across the 30 universes of a generator, the
values of $T$ for the two synthetic seeds have Spearman correlation between $-0.03$ and 0.30 (mean
0.14), and those of $\zhac$ between 0.10 and 0.51. For MMD$^2_u$ the values are 0.22 to 0.83 (mean
0.44). With the data held fixed (reference row), $T$ still has seed-to-seed correlation 0.28,
while MMD$^2_u$ is identical in both seeds (1.00): the randomness of the test alone, a split and a
network trained on few windows, makes $T$ barely reproducible. The link uses the mean of the two
seeds, which halves the variance of the seed noise but leaves it large relative to the variation
of $T$ across the universes of a generator (\cref{app:mk-link} quantifies the attenuation).

\begin{table}[tbp]
  \centering
  \caption{Spearman correlations between test quantities and utility losses (E4; sample
  moments). Panel A: replica-wise portfolios (the DESIGN measure); panel B: portfolios from the
  moments of the pooled replicas (misspecification only). Test quantities are seed means (two
  seeds); $\abs{\Delta\,\text{avg.\ corr.}}$: absolute difference between the synthetic and the
  real \TRAIN{} mean pairwise correlation (not a test). Pooled: over the 150 (universe,
  generator) points, universe-bootstrap 95\% percentile interval (2000 draws). Within-universe and
  within-generator: mean Spearman correlation across the 5 generators of a universe or the 30
  universes of a generator, two-sided permutation p-values (2000 permutations). Generator level:
  Spearman correlation of the 5 generator means, exact two-sided p-value over the 120 orderings.
  Universes share dates, so intervals are optimistic. All losses:
  \cref{tab:e4-link-full,tab:e4-link-pooled}.}
  \label{tab:e4-link}
  \small
  \fittable{manual_e4_link_main}
\end{table}

\paragraph{Correlations between test quantities and losses.}
\Cref{tab:e4-link} and \cref{fig:e4-utility-vs-distance} make the comparison systematic. For the
replica-wise GMV variance ratio, the \WC{} statistic $T$ has a pooled correlation of 0.10
[$-0.04$, 0.25], a within-generator correlation of 0.01 (permutation $p=0.853$) and a
generator-level correlation of 0.50 (exact $p=0.450$, $n=5$). Given the low reproducibility of
$T$, the within-generator value is weak evidence of no association: it is compatible with seed
noise alone and does not exclude a moderate association between the loss and a noise-free critic
distance. The MMD quantities, whose statistics are on average more reproducible, are no better:
both have within-generator correlations with the GMV variance ratio of $-0.08$ for replica-wise
and $-0.02$ for pooled-moment portfolios, and the wild-standardised MMD statistic is even
correlated negatively with the loss ($-0.21$ pooled). For the CE loss with $\gamma=5$ no test
quantity has a pooled correlation larger than 0.16 in absolute value. With pooled-moment
portfolios (panel B), $T$ orders the generators roughly as their GMV losses (pooled 0.30 [0.19,
0.42], within universe 0.39 with $p<0.001$, generator level 0.70 with exact $p=0.233$). Its
within-generator correlation, $-0.01$ ($p=0.895$), is weak evidence for the same reason, and the
between-generator association rests on five generators, of which only factor GARCH-$t$ has a
large loss. The evidence that test results do not predict utility thus comes from the comparison
of generators (MV Gaussian rejected everywhere yet with GMV variance ratios close to one, 1.038
replica-wise and 1.004 pooled; factor GARCH-$t$ rejected less often yet losing most) and from MMD, not from the within-generator correlations of the critic.
Over the full tables, 18 of the 192 combinations of test quantity, loss, aggregation and view have
a pooled interval excluding zero or a permutation p-value below 0.05. If all associations were null and
the comparisons independent, about eight would be expected by chance (\cref{app:mk-link}); the
comparisons share their points, so the count is descriptive. The four positive ones are all $T$
against pooled-moment GMV losses, in the pooled and within-universe views, which mostly rank the
generators; the other 14 are negative and all involve a studentised statistic ($\zhac$ or the
wild-standardised MMD statistic).

\begin{figure}[tbp]
  \centering
  \inputfigure{e4_utility_vs_distance}
  \caption{Utility losses against test statistics (E4), one point per (universe, generator), 30
  universes $\times$ 5 generators. Top: GMV variance ratio (TSTR over TRTR); bottom: mean-variance
  ($\gamma=5$) CE loss in percent per year; replica-wise portfolios, sample moments. Left: \WC{}
  statistic $T$; right: $\TMMD$; seed means over two seeds. Insets: pooled Spearman $\rho$ and
  within-generator $\bar\rho$ with universe-bootstrap 95\% percentile intervals. Dashed lines: no
  loss.}
  \label{fig:e4-utility-vs-distance}
\end{figure}

\paragraph{The correlation shift.}
The one quantity with a consistent association is not a test statistic.
$\abs{\Delta\,\text{avg.\ corr.}}$ is correlated with the GMV variance ratio at $-0.29$ [$-0.44$,
$-0.14$] pooled and $-0.36$ ($p<0.001$) within generator for replica-wise portfolios, and at
$-0.64$ [$-0.70$, $-0.57$] and $-0.41$ ($p<0.001$) for pooled-moment portfolios. The further a
generator's correlation is from the past, the smaller its GMV loss. The synthetic correlation lies
below the \TRAIN{} one in 122 of the 150 points (in all 30 universes for CCC-GARCH-$t$ and MV
Student-$t$), and the real correlation fell, so departing from the past meant moving towards the
future; a generator that reproduced the past correlation perfectly was penalised for it. The
association holds within generators, so it does not come from generator-level offsets, and it is
specific to GMV: for the CE loss with $\gamma=5$ it is $-0.03$ pooled. It is also a property of
this evaluation period, in which the correlation fell in each of the 30 universes; under rising
correlations the sign would plausibly reverse. The alternative split of \cref{sec:robust}, with
\TEST{} starting in 2022, was not run for E4.

\paragraph{Summary.}
Closeness in distribution does not imply usefulness for portfolio construction. The GMV portfolio
depends only on the covariance matrix up to scale, so features that a window test detects can be
irrelevant to it; mean-variance utility is driven by mean errors that no test sees at these sample
sizes; and when cross-sectional dependence shifts between calibration and evaluation, the
generator closest to the past is not the most useful. These findings are conditional on the
design. They come from one evaluation period, in which correlations fell, as point estimates
without inference over dates, and from five generators; in a period of rising correlations the
negative association between the correlation distance and the GMV loss could reverse. The tests
examine a \TEST{}-length sample started from the end-of-\TRAIN{} state, whereas the portfolios
use each generator only through the unconditional moments of long stationary samples, so the
conditional forecasts of the GARCH-type generators are not exploited. Within a generator, the multivariate critic statistic is so noisy
across seeds that its correlation with the loss is weak evidence either way, and the evidence
rests on the comparison of generators and on MMD. None of this is a failure of the test, which
answers, with an approximate size adjustment, whether synthetic windows are distinguishable from
real \TEST{} windows. Whether synthetic data are good for building portfolios has to be checked
on that task, by comparing TSTR with TRTR portfolios on held-out real data.
